\chapter*{Abstract}
\addcontentsline{toc}{chapter}{Abstract}

In this project, we present the design of a Retrieval-Augmented Generation (RAG) workflow integrated into a Visual Studio Code (VS Code) extension, serving as a local intelligent assistant for web developers working with the Next.js framework. The system uses Milvus \cite{wang2021milvus} as the vector database to efficiently store and retrieve semantically indexed information from the official Next.js documentation, prepared during the data processing phase.

For response generation, the architecture is model-agnostic and designed to work with locally hosted Large Language Models (LLMs) via Ollama\footnote{Ollama - \url{https://ollama.com/}}, a lightweight framework that allows users to run various LLMs depending on their hardware capabilities. This flexibility empowers developers to choose the model that best suits their system-whether it's a smaller model for lightweight environments or a more powerful one for advanced tasks.

By combining retrieval from up-to-date documentation with locally generated, context-aware responses, the proposed RAG-enhanced assistant boosts productivity and simplifies development workflows directly within the VS Code environment.

% Retrieval-Augmented Generation (RAG) has emerged as a transformative framework in the field of natural language processing, enabling large language models (LLMs) to leverage external knowledge effectively. This research explores the application of RAG techniques to enhance coding workflows within a specialized Visual Studio Code (VSCode) extension tailored for web development. By integrating state-of-the-art embedding models such as BGE-M3 and robust language generation models like Qwen, our system provides precise and context-aware assistance for developers. Key contributions include a modular architecture that combines vector databases, efficient indexing strategies, and RAG-based augmentation, along with a seamless integration of local LLMs through Ollama, ensuring data privacy and high performance. The research also addresses critical challenges, including managing rapidly evolving frameworks like Next.js and improving interoperability across tools such as Firebase and Tailwind CSS. Evaluation metrics focus on retrieval precision, response relevance, and computational efficiency, demonstrating the system’s potential to revolutionize development workflows. This study highlights the power of RAG in bridging the gap between generative AI and real-world coding applications, paving the way for more intelligent, adaptive, and developer-centric tools.